% ============================================================
%  Section 5 --- Results
%  Scaffold. Prose that does not depend on numbers is written;
%  every quantitative claim is a \TODO with the measurement named.
%  Target ~2 pages when filled.
% ============================================================

\section{Results}
\label{sec:results}

We report three things: whether the policy learns a placement rule at all, what
a single query costs it relative to two baselines, and what geometry the learned
rule produces. \Cref{sec:breaks} then reports where it stops working.

% ------------------------------------------------------------
\subsection{Setup}
\label{sec:results_setup}
% ------------------------------------------------------------

We solve the viscous Burgers equation $u_t + u\,u_x = \alpha u_{xx}$ on
$x \in [0,1]$ with $u(0,t) = u(1,t) = 0$ and a smooth initial condition, taking
the reference solution from the Hopf--Cole transform so that the error at a
query is measured against an analytical value rather than a discrete one.
\TODO{State $\alpha$, the initial condition, $T$, and the reference grid
$(n_x, \Delta x, \Delta t)$, with the CFL constraint that fixes $\Delta t$.}

What the agent is given, and what it must produce, should be stated precisely,
since it determines what the cost figures mean. At reset the visited set
contains the initial condition sampled on the reference grid and the boundary
values at every reference time. Everything in the interior is computed: each
interior point in $\mathcal{V}$ is the result of a reconstruction from points
already in $\mathcal{V}$, so errors propagate through the cloud rather than
being refreshed against ground truth. The reference solution is used only to
score the terminal prediction and never enters the rollout.

We compare against two baselines.

\emph{Masked-random} draws uniformly from the admissible actions at each step,
in the same environment, with the same masking and the same termination rule.
It isolates the contribution of the policy from the contribution of the
environment's constraints: any competence that survives when the policy is
replaced by a coin flip was never the policy's.

\emph{Uniform march} advances the same reconstruction operator on a regular
grid, which by \cref{sec:recon_pde} is the classical scheme. It is the
comparator that makes the cost figure meaningful in the units a numerical
analyst would use.

\TODO{Table: environment and training configuration for the featured run.
Resolve $R$, $K_{\max}$ and $e_{\mathrm{tol}}$ against the run actually
featured --- the notebook currently carries more than one value for each.}

% ------------------------------------------------------------
\subsection{Does the policy learn a placement rule?}
\label{sec:results_learn}
% ------------------------------------------------------------

We track three quantities over training: the fraction of episodes that
terminate with an admissible stencil at the query, the terminal error on those
episodes, and the number of evaluations consumed before termination. The first
measures whether the agent can construct a usable stencil at all; the second
whether it constructs a good one; the third what that costs.

\TODO{Figure: the three curves against environment steps, for the featured
horizon, with masked-random drawn as a horizontal band rather than a curve
(it does not improve). Caption should say which run.}

The comparison that matters is against masked-random rather than against
initialisation. Because every action advances in time and the mask already
removes collisions, out-of-domain proposals and cells whose solves have
recently failed, a policy that does nothing intelligent still terminates on a
large fraction of episodes. The question is not whether the reached rate is
high but whether it, and the error that accompanies it, are higher than what
the environment produces on its own.

\TODO{The headline comparison: reached rate and terminal error, learned policy
vs.\ masked-random, at matched budget. If the gap is small, say so plainly and
report it as a property of the task rather than of the method --- the
formulation makes reaching easy and accuracy hard, and that is itself worth
stating.}

\TODO{Policy entropy at convergence, against $\log(N|\mathcal{A}|)$. This is
the cheapest evidence that the policy has committed to something rather than
remaining near-uniform.}

% ------------------------------------------------------------
\subsection{The cost of a single query}
\label{sec:results_cost}
% ------------------------------------------------------------

The natural currency is the number of operator evaluations, since by
\cref{sec:mdp} that is what an action spends and what an episode is charged
for. We therefore report accuracy against evaluations rather than against wall
time, which depends on implementation details that are not the subject here.

\TODO{Figure: work--precision. Terminal error at $z^{*}$ against number of
evaluated points, three curves --- learned policy, masked-random, uniform march
--- with the uniform march swept over grid resolutions so that it traces a
curve rather than a point. Log--log axes. This is the headline claim of the
paper and the figure a reviewer will look at first.}

\TODO{One sentence quantifying it: to reach error $\varepsilon$ at $z^{*}$ the
learned policy evaluates $n_{\mathrm{RL}}$ points against $n_{\mathrm{unif}}$
for the uniform march, a factor of $\times$. State the horizon it holds at, and
whether the factor grows or shrinks with the horizon.}

% ------------------------------------------------------------
\subsection{What geometry does the policy discover?}
\label{sec:results_geometry}
% ------------------------------------------------------------

If the placement rule reflects the equation rather than the shortest route to
termination, the cloud it builds should occupy the region the solution at
$z^{*}$ actually depends on. That region has a known scale. For a
diffusion-dominated problem the numerically significant domain of dependence
widens as $\sqrt{\alpha\tau}$ with lookback $\tau = t^{*} - t$; for an
advection-dominated one it follows the characteristic and widens linearly at
the wave speed. The action set permits a far wider cone than either --- a walker
may move one cell in $x$ per time step, a speed of $\Delta x/\Delta t$ that
exceeds the physical scales by an order of magnitude --- so a narrower learned
envelope is informative rather than forced.

We therefore fit the half-width $w$ of the visited cloud against lookback as a
power law $w \propto \tau^{p}$ and read the exponent. Three outcomes are
distinguishable: $p \approx \tfrac12$ indicates the diffusive scale,
$p \approx 1$ with a slope matching the wave speed indicates characteristic
following, and $p \approx 1$ at the maximum action speed indicates that the
structure is an artefact of taking the shortest route to the query.

\TODO{Figure: visited cloud for a converged policy, with the fitted envelope
and the two physical scales overlaid; plus the log--log fit with $p$ and its
confidence interval.}

\TODO{Report $p$. If it comes out at the action-speed bound, report that --- it
is a negative result about the geometry, not about the method, and it is more
useful than silence.}

One caveat governs which horizon this can be measured at. At short horizons the
achieved error falls below $e_{\mathrm{tol}}$, so $e_{\mathrm{norm}} \equiv 0$
and the accuracy term in \eqref{eq:reward} contributes nothing: the agent is
optimising step count alone, and the compact cloud it produces is the shortest
route by construction. The fit is therefore only meaningful at horizons where
the accuracy term is active, and we report it there.

% ------------------------------------------------------------
\subsection{Ablations}
\label{sec:results_ablations}
% ------------------------------------------------------------

\TODO{One compact table rather than several plots. Rows worth including, all
of which exist in the training record:
\begin{itemize}
  \item \emph{Masking on/off} --- the livelock of \cref{sec:mdp_mask};
  \item \emph{State definition} --- with and without the query coordinates in
        the observation, which is a controlled pair at matched $N$, $\lambda$
        and $\beta$;
  \item \emph{Reward form} --- unnormalised versus normalised, where the
        unnormalised runs show the reached rate collapsing to zero as the
        accuracy weight grows;
  \item \emph{Bubble radius $R$} and \emph{walker count $N$} --- the two
        parameters that decide whether long horizons converge;
  \item \emph{Shaping coefficient} $c_{\mathrm{shape}}$, including
        $c_{\mathrm{shape}} = 0$.
\end{itemize}
Report reached rate and terminal error for each, at matched budget.}

Two of these rows are worth a sentence in the text rather than only a line in
the table. \TODO{Whichever of $R$ and $N$ turns out to separate converging from
non-converging runs, since \cref{sec:breaks} argues that this is a property of
the stencil geometry rather than of the learning problem.}
